\documentclass[11pt,a4paper]{article}
\usepackage[margin=1.8cm]{geometry}
\usepackage{amsmath,amssymb}
\usepackage{booktabs,array,tabularx}
\usepackage{tcolorbox}
\usepackage{microtype}
\usepackage{parskip}
\usepackage{enumitem}
\usepackage{xcolor}
\tcbuselibrary{skins,breakable}

% ── B&W only box styles ───────────────────────────────────────────────────────
\tcbset{
  defn/.style={colback=gray!8,colframe=black,fonttitle=\bfseries\small,
    title=#1,left=4pt,right=4pt,top=3pt,bottom=3pt,boxrule=0.7pt,arc=2pt},
  bold/.style={colback=white,colframe=black,fonttitle=\bfseries\small,
    title=#1,left=4pt,right=4pt,top=3pt,bottom=3pt,boxrule=1.2pt,arc=2pt},
  exam/.style={colback=gray!5,colframe=black,fonttitle=\bfseries\small,
    title=#1,left=4pt,right=4pt,top=3pt,bottom=3pt,
    boxrule=0.5pt,arc=2pt,borderline west={2pt}{0pt}{black}},
  note/.style={colback=white,colframe=gray!60,fonttitle=\bfseries\small,
    title=#1,left=4pt,right=4pt,top=3pt,bottom=3pt,boxrule=0.5pt,arc=2pt}
}
\setlist[itemize]{noitemsep,topsep=3pt,parsep=0pt}
\setlist[enumerate]{noitemsep,topsep=3pt,parsep=0pt}

\begin{document}

% ── Title ────────────────────────────────────────────────────────────────────
\begin{center}
  {\large\bfseries DRL Sessions 9 \& 10 — Temporal Difference Learning (Exam Handout)}\\[2pt]
  {\small S9--10: TD(0) · SARSA · Q-Learning · Maximization Bias · Race Car with TD}
\end{center}
\hrule\vspace{6pt}

% ════════════════════════════════════════════════════════════════════════════
\section*{1 · Importance of TD Methods in RL}
% ════════════════════════════════════════════════════════════════════════════

\begin{tcolorbox}[defn={Why TD is Central to RL}]
TD methods are the most important class of RL algorithms. They are:
\begin{itemize}
  \item \textbf{Model-free}: learn directly from experience — no $p(s',r|s,a)$ needed.
  \item \textbf{Online}: update at every step — no need to wait for episode end.
  \item \textbf{Bootstrapping}: use the current estimate of the next state's value to update
    the current state's value.
\end{itemize}
This combination makes TD practical for large, real-time problems where MC is too slow
(must wait for episode) and DP is infeasible (model unavailable).
\end{tcolorbox}

% ════════════════════════════════════════════════════════════════════════════
\section*{2 · How TD Connects to DP and MC}
% ════════════════════════════════════════════════════════════════════════════

\begin{tcolorbox}[bold={The Triangle: DP vs MC vs TD}]
\begin{center}
\begin{tabular}{p{3cm}p{4cm}p{4cm}p{3.5cm}}
\toprule
\textbf{Property} & \textbf{Dynamic Programming} & \textbf{Monte Carlo} & \textbf{TD} \\
\midrule
Needs model? & Yes & No & No \\
Bootstraps? & Yes & No & \textbf{Yes} \\
Online (step-by-step)? & Yes (sweep) & No (end of episode) & \textbf{Yes} \\
Target & $\mathbb{E}[R + \gamma V(S')]$ & Full return $G_t$ & $R_{t+1} + \gamma V(S_{t+1})$ \\
Bias & Low & None & Some (due to bootstrap) \\
Variance & Low & High & Low \\
\bottomrule
\end{tabular}
\end{center}
\medskip
\textbf{TD borrows from both:} uses experience like MC (model-free) and bootstraps like DP (one-step lookahead).
This is the key insight — it is ``online MC + bootstrapping''.
\end{tcolorbox}

\textbf{The update targets compared:}
\begin{align*}
  \text{DP target:}\quad      & \sum_{s',r} p(s',r|s,a)\bigl[r + \gamma V(s')\bigr] &&\text{(uses model, exact expectation)} \\
  \text{MC target:}\quad      & G_t = R_{t+1} + \gamma R_{t+2} + \cdots &&\text{(actual return, no bootstrap)} \\
  \text{TD(0) target:}\quad   & R_{t+1} + \gamma V(S_{t+1}) &&\text{(one sample, one bootstrap step)}
\end{align*}

% ════════════════════════════════════════════════════════════════════════════
\section*{3 · Framework of TD(0) — The Basic TD Method}
% ════════════════════════════════════════════════════════════════════════════

\begin{tcolorbox}[bold={TD(0) Prediction Update}]
\[
  V(S_t) \;\leftarrow\; V(S_t) + \alpha\underbrace{\bigl[R_{t+1} + \gamma V(S_{t+1}) - V(S_t)\bigr]}_{\delta_t \;=\; \text{TD error}}
\]
\begin{itemize}
  \item $\delta_t$ = \textbf{TD error} — how surprising the transition was (positive: better than expected, negative: worse).
  \item $R_{t+1} + \gamma V(S_{t+1})$ = the \textbf{TD target} (our improved estimate using next state).
  \item Update happens \emph{every step}, not at episode end.
  \item Converges to $v_\pi(s)$ as long as $\alpha$ decreases appropriately.
\end{itemize}
\end{tcolorbox}

\begin{tcolorbox}[exam={Worked Example — TD(0) update trace}]
\textbf{Given:} $V(S_t) = 0.5$, $R_{t+1} = 1$, $V(S_{t+1}) = 0.8$, $\alpha = 0.1$, $\gamma = 0.9$.

\[
  \delta_t = 1 + 0.9 \times 0.8 - 0.5 = 1 + 0.72 - 0.5 = 1.22
\]
\[
  V(S_t) \leftarrow 0.5 + 0.1 \times 1.22 = \mathbf{0.622}
\]
The state was better than expected (large positive reward + good next state) — value rises.

\medskip
\textbf{If} $R_{t+1} = -1$, $V(S_{t+1}) = 0.2$:
\[
  \delta_t = -1 + 0.9 \times 0.2 - 0.5 = -1 + 0.18 - 0.5 = -1.32
\]
\[
  V(S_t) \leftarrow 0.5 + 0.1 \times (-1.32) = \mathbf{0.368}
\]
Worse than expected — value falls.
\end{tcolorbox}

% ════════════════════════════════════════════════════════════════════════════
\section*{4 · SARSA — On-Policy TD Control}
% ════════════════════════════════════════════════════════════════════════════

\textbf{Goal:} Learn $q_\pi(s,a)$ (action-value function) to derive a control policy.
We need action-values (not state-values) since we have no model to compute $\arg\max_a$.

\begin{tcolorbox}[bold={SARSA Update Rule}]
After transition $(S_t, A_t) \xrightarrow{R_{t+1}} (S_{t+1}, A_{t+1})$:
\[
  Q(S_t, A_t) \;\leftarrow\; Q(S_t, A_t)
    + \alpha\bigl[R_{t+1} + \gamma Q(S_{t+1}, A_{t+1}) - Q(S_t, A_t)\bigr]
\]
\textbf{Name:} \textbf{S}tate, \textbf{A}ction, \textbf{R}eward, \textbf{S}tate, \textbf{A}ction
— the five elements of the update tuple.
\end{tcolorbox}

\textbf{Why on-policy?}\\
The next action $A_{t+1}$ is sampled from the \emph{current policy} $\pi$ (e.g.\ $\varepsilon$-greedy).
The update uses $Q(S_{t+1}, A_{t+1})$ — the value of the action \emph{actually taken next}.
So SARSA evaluates and improves the \emph{same} policy that generates experience.

\textbf{SARSA algorithm (pseudocode):}
\begin{enumerate}
  \item Initialise $Q(s,a)$ arbitrarily; $Q(\text{terminal},\cdot) = 0$.
  \item For each episode: initialise $S$; choose $A$ from $Q$ using $\varepsilon$-greedy.
  \item Repeat each step: take $A$, observe $R, S'$; choose $A'$ from $Q(S',\cdot)$ using $\varepsilon$-greedy.
  \item Apply SARSA update. Set $S \leftarrow S'$, $A \leftarrow A'$.
  \item Until $S$ is terminal.
\end{enumerate}

% ════════════════════════════════════════════════════════════════════════════
\section*{5 · Q-Learning — Off-Policy TD Control}
% ════════════════════════════════════════════════════════════════════════════

\begin{tcolorbox}[bold={Q-Learning Update Rule}]
After transition $(S_t, A_t) \xrightarrow{R_{t+1}} S_{t+1}$:
\[
  Q(S_t, A_t) \;\leftarrow\; Q(S_t, A_t)
    + \alpha\bigl[R_{t+1} + \gamma \max_{a} Q(S_{t+1}, a) - Q(S_t, A_t)\bigr]
\]
\textbf{Key difference from SARSA:} uses $\max_a Q(S_{t+1}, a)$, not $Q(S_{t+1}, A_{t+1})$.
It always bootstraps from the \emph{best} possible next action, regardless of what the
behaviour policy actually does.
\end{tcolorbox}

\textbf{Why off-policy?}\\
The target $\max_a Q(S_{t+1}, a)$ is the Bellman optimality equation — it assumes the agent
will act greedily from the next state. But the actual next action $A_{t+1}$ may be exploratory.
Q-Learning learns $q^*$ \emph{directly}, independently of the policy being followed.
The behaviour policy just needs to cover all $(s,a)$ pairs.

% ════════════════════════════════════════════════════════════════════════════
\section*{6 · SARSA vs.\ Q-Learning — Key Differences}
% ════════════════════════════════════════════════════════════════════════════

\begin{tcolorbox}[defn={Side-by-Side Comparison}]
\begin{tabular}{p{3.8cm}p{4.8cm}p{5cm}}
\toprule
\textbf{Property} & \textbf{SARSA} & \textbf{Q-Learning} \\
\midrule
Update uses & $Q(S_{t+1}, A_{t+1})$ — actual next action & $\max_a Q(S_{t+1}, a)$ — best possible action \\
Policy type & On-policy & Off-policy \\
What it learns & $q_\pi$ (value of \emph{followed} policy) & $q^*$ (value of \emph{optimal} policy) \\
Behaviour on Cliff Walk & Takes safer inland path (cautious) & Falls off cliff under $\varepsilon$-greedy (learns optimal but risky path) \\
Convergence & Converges to $q^*$ under GLIE & Converges to $q^*$ directly \\
Safe / cautious? & Yes — accounts for actual exploration & No — ignores exploration risk \\
\bottomrule
\end{tabular}
\end{tcolorbox}

\begin{tcolorbox}[exam={Cliff Walking Example — SARSA vs Q-Learning (must know)}]
\textbf{Setup:} Grid world. Bottom row has a cliff. Agent starts at bottom-left, goal is bottom-right.
Falling off cliff: reward $-100$, restart. All other moves: reward $-1$.

\medskip
\textbf{Q-Learning} learns the \emph{theoretically} optimal path — walk along the cliff edge
(fewest steps, reward $\approx -13$). But under $\varepsilon$-greedy, the agent occasionally steps
off the cliff (getting $-100$). \emph{Online performance is poor.}

\medskip
\textbf{SARSA} accounts for the exploration policy. It ``knows'' the agent sometimes makes random moves,
so it learns to avoid the cliff edge — takes the safer inland path (reward $\approx -17$ but stays
safe during learning). \emph{Online performance is better.}

\medskip
\textbf{Exam takeaway:}
SARSA $=$ safe during training (on-policy). Q-Learning $=$ finds optimal but risky path during training.
As $\varepsilon \to 0$, both converge to the same optimal policy.
\end{tcolorbox}

\begin{tcolorbox}[exam={Worked Trace — SARSA vs Q-Learning on same transition}]
\textbf{Given:} $Q(S_t, A_t) = 2.0$, $\alpha = 0.1$, $\gamma = 0.9$, $R_{t+1} = 1$.\\
Next state $S_{t+1}$ has $Q(S_{t+1}, \text{left}) = 3.0$, $Q(S_{t+1}, \text{right}) = 5.0$.\\
Actual next action (by $\varepsilon$-greedy): $A_{t+1} = \text{left}$ (exploratory choice).

\medskip
\textbf{SARSA} (uses actual $A_{t+1} = \text{left}$):
\[
  Q(S_t, A_t) \leftarrow 2.0 + 0.1[1 + 0.9 \times 3.0 - 2.0]
             = 2.0 + 0.1[1 + 2.7 - 2.0]
             = 2.0 + 0.1 \times 1.7 = \mathbf{2.17}
\]

\textbf{Q-Learning} (uses $\max = 5.0$):
\[
  Q(S_t, A_t) \leftarrow 2.0 + 0.1[1 + 0.9 \times 5.0 - 2.0]
             = 2.0 + 0.1[1 + 4.5 - 2.0]
             = 2.0 + 0.1 \times 3.5 = \mathbf{2.35}
\]

Q-Learning gives a higher estimate (more optimistic) because it uses the \emph{best} possible
next action, not the one actually taken.
\end{tcolorbox}

% ════════════════════════════════════════════════════════════════════════════
\section*{7 · Maximization Bias and Double Q-Learning}
% ════════════════════════════════════════════════════════════════════════════

\subsection*{7.1 · The Problem: Maximization Bias}

\begin{tcolorbox}[defn={What is Maximization Bias?}]
Q-Learning uses $\max_a Q(S_{t+1}, a)$ in its update. But $Q(s,a)$ is a \emph{noisy estimate}
of $q^*(s,a)$.

\medskip
\textbf{The bias:}
\[
  \mathbb{E}\bigl[\max_a Q(s,a)\bigr] \;\geq\; \max_a\,\mathbb{E}[Q(s,a)] = \max_a q^*(s,a)
\]
The expected maximum of noisy estimates is \textbf{always} $\geq$ the true maximum.

\medskip
\textbf{Why?} Even if all actions have the same true value, random noise will make some look
better. Selecting the maximum \emph{amplifies} that noise — we systematically pick the
overestimated action.
\end{tcolorbox}

\textbf{Impact on learning:}
\begin{itemize}
  \item Agent is over-optimistic about values of certain states/actions.
  \item Can cause the agent to commit to suboptimal actions because they \emph{look} good early on.
  \item Performance degrades or converges more slowly.
  \item More severe when there are many actions with similar true values.
\end{itemize}

\begin{tcolorbox}[exam={Illustrative Example — Maximization Bias}]
\textbf{Setup:} One state $A$, two actions: \texttt{left} (goes to state $B$) and \texttt{right}
(terminal, reward $= 0$). State $B$ has many actions, each with reward drawn from
$\mathcal{N}(-0.1, 1)$ (mean $= -0.1$, noisy).

\medskip
\textbf{True value:} $Q(A, \text{left}) \approx -0.1$ (negative in expectation).
$Q(A, \text{right}) = 0$.

\textbf{What Q-Learning does early:} With few samples, some actions in $B$ appear positive
(lucky noise). So $\max_a Q(B, a) > 0$, making $Q(A, \text{left}) > 0$.
Agent incorrectly prefers \texttt{left} over \texttt{right}.

\textbf{Maximization bias:} Q-Learning overestimates $Q(A, \text{left})$ due to noise.
\end{tcolorbox}

\subsection*{7.2 · Solution: Double Q-Learning}

\begin{tcolorbox}[bold={Double Q-Learning — Separating Selection and Evaluation}]
Maintain two independent value estimates $Q_1$ and $Q_2$.

\medskip
At each step, with probability 0.5:
\[
  Q_1(S_t, A_t) \leftarrow Q_1(S_t, A_t)
    + \alpha\Bigl[R_{t+1} + \gamma Q_2\!\bigl(S_{t+1},\, \arg\max_a Q_1(S_{t+1}, a)\bigr)
      - Q_1(S_t, A_t)\Bigr]
\]
Otherwise swap $Q_1 \leftrightarrow Q_2$ in the update.

\medskip
\textbf{Key idea:} $Q_1$ selects the action ($\arg\max_a Q_1$); $Q_2$ evaluates it.
Since $Q_1$ and $Q_2$ are independent, the bias is removed:
\[
  \mathbb{E}\bigl[Q_2(s, \arg\max_a Q_1(s,a))\bigr] = q^*(s, \arg\max_a q^*(s,a))
\]
No systematic overestimation.
\end{tcolorbox}

% ════════════════════════════════════════════════════════════════════════════
\section*{8 · Race Car Example with TD}
% ════════════════════════════════════════════════════════════════════════════

\begin{tcolorbox}[defn={Race Track with SARSA and Q-Learning}]
\textbf{State:} $(x, y, v_x, v_y)$ — position + velocity.
\textbf{Actions:} 9 acceleration combinations.
\textbf{Reward:} $-1$ per step; going off-track: $-1$ + return to start.
\end{tcolorbox}

\textbf{TD vs.\ MC on the race track:}
\begin{center}\small
\begin{tabular}{p{3cm}p{5.5cm}p{5.5cm}}
\toprule
& \textbf{MC (from Session 8)} & \textbf{TD (this session)} \\
\midrule
When it updates & End of each lap & Every step (online) \\
Episodes needed & Many full laps required & Learns even from partial runs \\
Handling crashes & Episode ends; returns -1 per step & Updates immediately on crash reward \\
Practical speed & Slow for long tracks & Much faster convergence \\
\bottomrule
\end{tabular}
\end{center}

\textbf{SARSA on race track:}\\
Agent follows $\varepsilon$-greedy policy. After each step, updates $Q(s,a)$ using actual next action.
SARSA learns to \emph{decelerate before corners} — it accounts for exploration errors (sometimes
veering off-track) and prefers safer trajectories.

\textbf{Q-Learning on race track:}\\
Uses $\max_a Q(S_{t+1}, a)$ at each step. Learns the theoretically optimal racing line (fastest path).
However, during learning under $\varepsilon$-greedy, it may crash more often on corners because it
ignores the exploration risk — analogous to the Cliff Walking result.

\begin{tcolorbox}[exam={Simplified 1-D Race Car Trace — Q-Learning}]
\textbf{States:} positions 0--6, velocity $v \in \{1,2\}$. Actions: \texttt{Acc} or \texttt{Hold}.
\textbf{Reward:} $-1$/step. $\alpha=0.1$, $\gamma=0.9$.
Initialise all $Q(s,a)=0$.

\medskip
\textbf{Step 1:} In state $(0,1)$, take \texttt{Acc}. Move to $(2,2)$. Reward $=-1$.\\
$\max_a Q((2,2), a) = 0$.
\[
  Q((0,1),\texttt{Acc}) \leftarrow 0 + 0.1[-1 + 0.9 \times 0 - 0] = -0.1
\]

\textbf{Step 2:} In state $(2,2)$, take \texttt{Hold}. Move to $(4,2)$. Reward $=-1$.\\
$\max_a Q((4,2), a) = 0$.
\[
  Q((2,2),\texttt{Hold}) \leftarrow 0 + 0.1[-1 + 0.9 \times 0 - 0] = -0.1
\]

\textbf{Step 3:} In state $(4,2)$, take \texttt{Hold}. Reach finish $(6,2)$. Reward $=-1$.
Terminal: $\max Q(\text{finish}) = 0$.
\[
  Q((4,2),\texttt{Hold}) \leftarrow 0 + 0.1[-1 + 0 - 0] = -0.1
\]

\medskip
After this first episode: values $= -0.1$ along the path. After many episodes, values become
increasingly negative for states far from finish (more steps $\Rightarrow$ lower return).
Greedy policy: \texttt{Acc} whenever possible (fewer steps $\Rightarrow$ higher value $\approx -$fewer steps).
\end{tcolorbox}

% ════════════════════════════════════════════════════════════════════════════
\section*{9 · Exam Key-Takeaways — Sessions 9 \& 10}
% ════════════════════════════════════════════════════════════════════════════

\begin{tcolorbox}[exam={Must Know — TD Methods}]
\begin{enumerate}
  \item \textbf{TD = model-free + bootstrapping + online.} Combines the best of MC and DP.
  \item \textbf{TD(0) update:} $V(S_t) \leftarrow V(S_t) + \alpha[\underbrace{R_{t+1} + \gamma V(S_{t+1})}_{\text{TD target}} - V(S_t)]$.
  \item \textbf{TD error} $\delta_t = R_{t+1} + \gamma V(S_{t+1}) - V(S_t)$ — know this symbol.
  \item \textbf{SARSA} (on-policy): update uses actual next action $A_{t+1}$.
    $Q(S,A) \leftarrow Q(S,A) + \alpha[R + \gamma Q(S',A') - Q(S,A)]$.
  \item \textbf{Q-Learning} (off-policy): update uses $\max_a Q(S',a)$, not actual $A'$.
    Learns $q^*$ directly regardless of behaviour policy.
  \item \textbf{Cliff Walk:} SARSA = safe inland path; Q-Learning = risky optimal path.
    SARSA has better \emph{online} performance under $\varepsilon$-greedy.
  \item \textbf{Maximization bias:} $\mathbb{E}[\max_a Q(s,a)] \geq \max_a q^*(s,a)$.
    Q-Learning is over-optimistic; hurts convergence and safety.
  \item \textbf{Double Q-Learning:} two independent estimators — one selects action, other evaluates it. Removes bias.
  \item \textbf{Race track with SARSA:} conservative (decelerate near corners);
    Q-Learning: finds fastest line but more crashes during training.
  \item \textbf{Key difference SARSA vs Q-Learning} in one word: SARSA is \emph{cautious};
    Q-Learning is \emph{optimistic}.
\end{enumerate}
\end{tcolorbox}

\vspace{4pt}
\hrule
\vspace{3pt}
{\footnotesize DRL EC3 Exam Handout — Sessions 9 \& 10 (Temporal Difference Learning). Ref: Sutton \& Barto Ch.~6.}

\end{document}
